\documentclass[10pt,a4paper]{amsart}
\usepackage{amsmath,amssymb,amsthm,graphicx}
\usepackage{verbatim}
\newenvironment{reference}{\comment}{\endcomment}
%\newenvironment{reference}{}{}
\newenvironment{slogan}{\comment}{\endcomment}
\newenvironment{history}{\comment}{\endcomment}

% For commutative diagrams we use Xy-pic
\usepackage[all]{xy}

% We use 2cell for 2-commutative diagrams.
\xyoption{2cell}
% Theorem environments.
%
\theoremstyle{plain}
\newtheorem{theorem}[subsection]{Theorem}
\newtheorem{proposition}[subsection]{Proposition}
\newtheorem{lemma}[subsection]{Lemma}

\theoremstyle{definition}
\newtheorem{definition}[subsection]{Definition}
\newtheorem{example}[subsection]{Example}
\newtheorem{exercise}[subsection]{Exercise}
\newtheorem{situation}[subsection]{Situation}

\theoremstyle{remark}
\newtheorem{remark}[subsection]{Remark}
\newtheorem{remarks}[subsection]{Remarks}

\numberwithin{equation}{subsection}

% Macros
%
\def\lim{\mathop{\mathrm{lim}}\nolimits}
\def\colim{\mathop{\mathrm{colim}}\nolimits}
\def\Spec{\mathop{\mathrm{Spec}}}
\def\Hom{\mathop{\mathrm{Hom}}\nolimits}
\def\Ext{\mathop{\mathrm{Ext}}\nolimits}
\def\SheafHom{\mathop{\mathcal{H}\!\mathit{om}}\nolimits}
\def\SheafExt{\mathop{\mathcal{E}\!\mathit{xt}}\nolimits}
\def\Sch{\mathit{Sch}}
\def\Mor{\mathop{\mathrm{Mor}}\nolimits}
\def\Ob{\mathop{\mathrm{Ob}}\nolimits}
\def\Sh{\mathop{\mathit{Sh}}\nolimits}
\def\NL{\mathop{N\!L}\nolimits}
\def\CH{\mathop{\mathrm{CH}}\nolimits}
\def\proetale{{pro\text{-}\acute{e}tale}}
\def\etale{{\acute{e}tale}}
\def\QCoh{\mathit{QCoh}}
\def\Ker{\mathop{\mathrm{Ker}}}
\def\Im{\mathop{\mathrm{Im}}}
\def\Coker{\mathop{\mathrm{Coker}}}
\def\Coim{\mathop{\mathrm{Coim}}}

% Boxtimes
%
\DeclareMathSymbol{\boxtimes}{\mathbin}{AMSa}{"02}

%
% Macros for moduli stacks/spaces
%
\def\QCohstack{\mathcal{QC}\!\mathit{oh}}
\def\Cohstack{\mathcal{C}\!\mathit{oh}}
\def\Spacesstack{\mathcal{S}\!\mathit{paces}}
\def\Quotfunctor{\mathrm{Quot}}
\def\Hilbfunctor{\mathrm{Hilb}}
\def\Curvesstack{\mathcal{C}\!\mathit{urves}}
\def\Polarizedstack{\mathcal{P}\!\mathit{olarized}}
\def\Complexesstack{\mathcal{C}\!\mathit{omplexes}}
% \Pic is the operator that assigns to X its picard group, usage \Pic(X)
% \Picardstack_{X/B} denotes the Picard stack of X over B
% \Picardfunctor_{X/B} denotes the Picard functor of X over B
\def\Pic{\mathop{\mathrm{Pic}}\nolimits}
\def\Picardstack{\mathcal{P}\!\mathit{ic}}
\def\Picardfunctor{\mathrm{Pic}}
\def\Deformationcategory{\mathcal{D}\!\mathit{ef}}
\usepackage{cleveref}
\crefname{originalsection}{original section}{original sections}
\Crefname{originalsection}{Original section}{Original sections}
\makeatletter
\newlabel{formal-spaces-example-reduction-affine-formal-spectrum}{{0GB6}{source}{}{sourcepointer-1}{}}
\newlabel{formal-spaces-example-reduction-affine-formal-spectrum@cref}{{[originalsection][0][]0GB6}{[1][0][]source}{}{}{}}
\makeatother
\crefname{originallemma}{original lemma}{original lemmas}
\Crefname{originallemma}{Original lemma}{Original lemmas}
\makeatletter
\newlabel{formal-spaces-lemma-etale}{{0AN8}{source}{}{sourcepointer-2}{}}
\newlabel{formal-spaces-lemma-etale@cref}{{[originallemma][0][]0AN8}{[1][0][]source}{}{}{}}
\makeatother
\crefname{originallemma}{original lemma}{original lemmas}
\Crefname{originallemma}{Original lemma}{Original lemmas}
\makeatletter
\newlabel{formal-spaces-lemma-etale-surjective}{{0AN9}{source}{}{sourcepointer-3}{}}
\newlabel{formal-spaces-lemma-etale-surjective@cref}{{[originallemma][0][]0AN9}{[1][0][]source}{}{}{}}
\makeatother
\crefname{originallemma}{original lemma}{original lemmas}
\Crefname{originallemma}{Original lemma}{Original lemmas}
\makeatletter
\newlabel{formal-spaces-lemma-reduction-smooth}{{0GB7}{source}{}{sourcepointer-4}{}}
\newlabel{formal-spaces-lemma-reduction-smooth@cref}{{[originallemma][0][]0GB7}{[1][0][]source}{}{}{}}
\makeatother
\crefname{originalreference}{original reference}{original references}
\Crefname{originalreference}{Original reference}{Original references}
\makeatletter
\newlabel{formal-spaces-situation-local-property}{{0CBA}{source}{}{sourcepointer-5}{}}
\newlabel{formal-spaces-situation-local-property@cref}{{[originalreference][0][]0CBA}{[1][0][]source}{}{}{}}
\makeatother
\crefname{originalsection}{original section}{original sections}
\Crefname{originalsection}{Original section}{Original sections}
\makeatletter
\newlabel{spaces-morphisms-section-finite-type}{{03XE}{source}{}{sourcepointer-6}{}}
\newlabel{spaces-morphisms-section-finite-type@cref}{{[originalsection][0][]03XE}{[1][0][]source}{}{}{}}
\makeatother
\title{Finite type on reduced formal spectra is a local property}
\author{The Stacks Project Authors}
\date{stacks-project:restricted.tex:local-snapshot-20261012; GNU FDL 1.2 or later}
\begin{document}
\maketitle
\noindent\textit{Editorial note.} All delivered mathematical text is the author's own, copied from the pinned Stacks Project chapter file restricted.tex; no statement or proof step was rewritten, reordered or removed.\par\medskip
\noindent\textit{Editorial note.} The counted proof is the author's complete argument: the proof is complete for all three clauses\par\medskip
\noindent\textit{Editorial note.} The label lemma-finite-type-red names the counted statement, and the proof environment immediately following it in the source is delivered verbatim as its proof.\par\medskip
\noindent\textit{Editorial note.} The counted statement is carried by the author's own lemma environment in the chapter source.\par\medskip
\noindent\textit{Allowed standard prerequisites (editorial).}\begin{itemize}
\item the WAdm\textasciicircum{}count category and the ideals of topologically nilpotent elements
\item formal spectra and their reduced closed subschemes
\item etale charts of affine formal algebraic spaces and etaleness of the induced maps on reductions
\item locality on source and target in the etale topology for locally finite type morphisms
\end{itemize}\medskip
\ifdefined\hypertarget\hypertarget{sourcepointer-1}{}\fi
\noindent\textit{Source reference.} Original section 0GB6: sources/stacks-chow/tags.txt, source lines 20361--20361. Prerequisite pointer invoked at proof line(s) 2676; it names a standard fact of a sibling chapter that the counted argument uses, and every step of the counted proof built on it is retained verbatim in the delivered body.\par\medskip
\ifdefined\hypertarget\hypertarget{sourcepointer-2}{}\fi
\noindent\textit{Source reference.} Original lemma 0AN8: sources/stacks-chow/tags.txt, source lines 13431--13431. Prerequisite pointer invoked at proof line(s) 2663; it names a standard fact of a sibling chapter that the counted argument uses, and every step of the counted proof built on it is retained verbatim in the delivered body.\par\medskip
\ifdefined\hypertarget\hypertarget{sourcepointer-3}{}\fi
\noindent\textit{Source reference.} Original lemma 0AN9: sources/stacks-chow/tags.txt, source lines 13432--13432. Prerequisite pointer invoked at proof line(s) 2677; it names a standard fact of a sibling chapter that the counted argument uses, and every step of the counted proof built on it is retained verbatim in the delivered body.\par\medskip
\ifdefined\hypertarget\hypertarget{sourcepointer-4}{}\fi
\noindent\textit{Source reference.} Original lemma 0GB7: sources/stacks-chow/tags.txt, source lines 20362--20362. Prerequisite pointer invoked at proof line(s) 2674; it names a standard fact of a sibling chapter that the counted argument uses, and every step of the counted proof built on it is retained verbatim in the delivered body.\par\medskip
\ifdefined\hypertarget\hypertarget{sourcepointer-5}{}\fi
\noindent\textit{Source reference.} Original reference 0CBA: sources/stacks-chow/tags.txt, source lines 15465--15465. Prerequisite pointer invoked at proof line(s) 2639; it names a standard fact of a sibling chapter that the counted argument uses, and every step of the counted proof built on it is retained verbatim in the delivered body.\par\medskip
\ifdefined\hypertarget\hypertarget{sourcepointer-6}{}\fi
\noindent\textit{Source reference.} Original section 03XE: sources/stacks-chow/tags.txt, source lines 5145--5145. Prerequisite pointer invoked at proof line(s) 2692; it names a standard fact of a sibling chapter that the counted argument uses, and every step of the counted proof built on it is retained verbatim in the delivered body.\par\medskip
\begin{lemma}
\label{lemma-finite-type-red}
For an arrow $\varphi : A \to B$ in $\text{WAdm}^{count}$ consider
the property $P(\varphi)=$``the induced ring homomorphism
$A/\mathfrak a \to B/\mathfrak b$ is of finite type''
where $\mathfrak a \subset A$ and $\mathfrak b \subset B$ are the ideals
of topologically nilpotent elements. Then $P$ is a local property
as defined in
Formal Spaces, Situation
\ref{formal-spaces-situation-local-property}.
\end{lemma}

\begin{proof}
Consider a commutative diagram
$$
\xymatrix{
B \ar[r] & (B')^\wedge \\
A \ar[r] \ar[u]^\varphi & (A')^\wedge \ar[u]_{\varphi'}
}
$$
as in Formal Spaces, Situation
\ref{formal-spaces-situation-local-property}.
Taking $\text{Spf}$ of this diagram we obtain
$$
\xymatrix{
\text{Spf}(B) \ar[d] &
\text{Spf}((B')^\wedge) \ar[l] \ar[d] \\
\text{Spf}(A) &
\text{Spf}((A')^\wedge) \ar[l]
}
$$
of affine formal algebraic spaces whose horizontal arrows are
representable by algebraic spaces and \'etale by
Formal Spaces, Lemma \ref{formal-spaces-lemma-etale}.
Hence we obtain a commutative diagram of affine schemes
$$
\xymatrix{
\text{Spf}(B)_{red} \ar[d]^f &
\text{Spf}((B')^\wedge)_{red} \ar[l]^g \ar[d]^{f'} \\
\text{Spf}(A)_{red} &
\text{Spf}((A')^\wedge)_{red} \ar[l]
}
$$
whose horizontal arrows are \'etale by
Formal Spaces, Lemma \ref{formal-spaces-lemma-reduction-smooth}.
By Formal Spaces, Example
\ref{formal-spaces-example-reduction-affine-formal-spectrum} and
Lemma \ref{formal-spaces-lemma-etale-surjective}
conditions (1), (2), and (3) of
Formal Spaces, Situation
\ref{formal-spaces-situation-local-property}
translate into the following statements
\begin{enumerate}
\item if $f$ is locally of finite type, then $f'$ is locally of finite type,
\item if $f'$ is locally of finite type and $g$ is surjective, then
$f$ is locally of finite type, and
\item if $T_i \to S$, $i = 1, \ldots, n$ are locally of finite type, then
$\coprod_{i = 1, \ldots, n} T_i \to S$ is locally of finite type.
\end{enumerate}
Properties (1) and (2) follow from the fact that being locally
of finite type is local on the source and target in the
\'etale topology, see discussion in
Morphisms of Spaces, Section \ref{spaces-morphisms-section-finite-type}.
Property (3) is a straightforward consequence of the definition.
\end{proof}

\end{document}
